\section{Agent-First 视角下的问题重构}

\subsection{从模型中心到系统中心}
课程提出 ``Agent-first'' 视角：将智能体视作一个持续运行的决策系统，而非一次性问答接口。该系统至少包含五个基本要素：状态表示、动作空间、记忆机制、规划机制与外部反馈回路。只有这五个要素协同，Agent 才能跨越短任务到长任务的鸿沟。

\begin{knowledgebox}{Agent 的最小系统抽象}
可以将语言智能体简化为如下循环：
\[
\text{Observe} \rightarrow \text{Reason} \rightarrow \text{Act} \rightarrow \text{Update Memory} \rightarrow \text{Replan}
\]
其中每个环节都可被独立增强，但系统性能取决于各环节之间的接口质量和延迟成本。
\end{knowledgebox}

\subsection{历史脉络与范式演进}
从早期 AutoGPT 到当前结合 ReAct、Tool Use、Self-Reflection 的新系统，范式演进并不是 ``算法替代算法''，更像 ``能力叠加能力''。课程强调三种增长路径：推理质量增长、记忆可靠性增长、规划尺度增长。它们共同决定可解决任务的上限。

\begin{table}[H]
\centering
\begin{tabular}{p{2.5cm}p{4.2cm}p{4.7cm}}
\toprule
\textbf{阶段} & \textbf{代表能力} & \textbf{主要瓶颈} \\
\midrule
早期 Agent 原型 & Prompt + 工具调用 & 长链任务中易漂移、易死循环 \\
中期工程化阶段 & ReAct + Reflection + Retrieval & 状态污染、成本高、调参复杂 \\
当前系统化阶段 & Reasoning actions + memory architecture + replanning & 验证器质量与跨场景泛化仍不足 \\
\bottomrule
\end{tabular}
\caption{Language Agent 的工程演进，不是单点突破，而是系统要素同步升级}
\end{table}

\begin{warningbox}{常见误解：把 Agent 当成“会调用工具的聊天机器人”}
这类定义忽略了长期状态管理和规划反馈。真正的 Agent 系统必须显式处理失败恢复、策略切换、记忆更新和目标重设，否则能力会停留在短任务。
\end{warningbox}

\subsection{本章小结}
Agent-first 视角把问题从 ``模型会不会答'' 转成 ``系统能否稳定完成任务''。这为后续三大能力模块提供了统一评估口径。

